\scanpage{20}
\mspage{--- 48 ---}
1) Le groupe $\mathrm{G}_{\mathrm{V}}$ est un sous-groupe \textit{ouvert} du groupe $\mathrm{H}_{\mathrm{V}}(\mathbf{Q}_p)$ des points de $\mathrm{H}_{\mathrm{V}}$ à coefficients dans $\mathbf{Q}_p$.

Il revient au même de dire que l’algèbre de Lie du groupe analytique $p$-adique $\mathrm{G}_{\mathrm{V}}$ est \textit{algébrique}, au sens de Chevalley.

2) La décomposition $\mathrm{V}_{\mathrm{C}}=\sum\mathrm{V}_{\mathrm{C}}(n)$ définit un homomorphisme
\[
h:\mathbf{G}_m/_{\mathrm{C}}\longrightarrow\mathbf{GL}_{\mathrm{V}}/_{\mathrm{C}}
\]
de groupes algébriques sur $\mathrm{C}$. Soit $\mathrm{H}_{\mathrm{V}}^{\unsure{\mathrm{O}|0|\circ}}$ la composante neutre de $\mathrm{H}_{\mathrm{V}}$, et soit $\mathrm{H}_{\mathrm{V}}^{\mathrm{O}}/_{\mathrm{C}}$ le sous-groupe de $\mathbf{GL}_{\mathrm{V}}/_{\mathrm{C}}$ déduit de $\mathrm{H}_{\mathrm{V}}^{\mathrm{O}}$ par extension des scalaires de $\mathbf{Q}_p$ à $\mathrm{C}$. Alors $\mathrm{H}_{\mathrm{V}}^{\mathrm{O}}/_{\mathrm{C}}$ \textit{contient le groupe} $h(\mathbf{G}_m/_{\mathrm{C}})$, image de $h$.

Lorsque de plus $\mathrm{G}_{\mathrm{V}}$ est \textit{résoluble}, on peut caractériser $\mathrm{H}_{\mathrm{V}}^{\mathrm{O}}$ comme \textit{le plus petit sous-groupe algébrique} $\mathrm{R}$ sur $\mathbf{Q}_p$ de $\mathbf{GL}_{\mathrm{V}}$ \textit{tel que} $\mathrm{R}/_{\mathrm{C}}$ \textit{contienne} $h(\mathbf{G}_m/_{\mathrm{C}})$ ; c’est l’analogue $p$-adique du \textit{groupe de Mumford-Tate}.

Il est probable que cette caractérisation de $\mathrm{H}_{\mathrm{V}}^{\mathrm{O}}$ vaut sans supposer $\mathrm{G}_{\mathrm{V}}$ résoluble ; on pourrait le prouver si l’on savait que $\mathrm{V}_{\mathrm{C}}=\mathrm{V}_{\mathrm{C}}(0)$ entraîne que $\mathrm{G}_{\mathrm{V}}$ est fini.
% slanted
\reversemarginpar
\margin{\handwritten{\shortstack[l]{ceci a été\\démontré par\\S. Sen (Annals)}\,$\left[\vphantom{\begin{matrix}1\\1\\1\end{matrix}}\right.$}}
\normalmarginpar

3) Lorsque $\mathrm{V}$ varie, les groupes $\mathrm{H}_{\mathrm{V}}$ forment un système projectif ; leur limite projective est un \textit{groupe proalgébrique} $\mathrm{H}$ sur $\mathbf{Q}_p$, muni d’un homomorphisme $h:\mathbf{G}_m/_{\mathrm{C}}\longrightarrow\mathrm{H}/_{\mathrm{C}}$. La catégorie des modules de Hodge est équivalente à celle des $\mathrm{H}$-modules de dimension finie ; si $\mathrm{H}^{\unsure{\mathrm{o}|\circ}}$ désigne la composante neutre de $\mathrm{H}$, le groupe $\mathrm{H}/\mathrm{H}^{\mathrm{o}}$ est le groupe de dimension zéro associé au groupe profini $\mathrm{G}$ ; le groupe $\mathrm{H}^{\mathrm{o}}$ ne change pas lorsqu’on remplace $\mathrm{K}$ par une extension finie.

4) Soit $\Theta$ le plus grand quotient commutatif du groupe proalgébrique $\mathrm{H}^{\mathrm{o}}$. Le groupe $\Theta$ est une limite projective de tores ; il ne dépend que de $p$ (mais pas de $\mathrm{K}$) ; on a :
\[
\Theta=\underset{\leftarrow}{\mathrm{lim.}}\,\mathrm{T}_{\mathrm{E}}
\]
où $\mathrm{E}$ parcourt l’ensemble des extensions finies de $\mathbf{Q}_p$.

(Si $\mathrm{E}$ est une extension de $\mathbf{Q}_p$ de degré $n$, on note $\mathrm{T}_{\mathrm{E}}$ le tore $\mathrm{R}_{\mathrm{E}}/_{\mathbf{Q}_p}(\mathbf{G}_m/_{\mathrm{E}})$ défini par $\mathrm{E}$ ; sa dimension est $n$, et $\mathrm{T}_{\mathrm{E}}(\mathbf{Q}_p)$ est canoniquement isomorphe à $\mathrm{E}^{*}$.)

Ce résultat est essentiellement équivalent à un théorème exposé par \textsc{Tate} dans le séminaire 1965/66. Il est intimement lié aux propriétés des « multiplications complexes formelles » des groupes $p$-divisibles.
% archive sheet number: 19
